\chapter{Protocole d'échange}
\label{ann:protocole}

\section{Commandes WebSocket (navigateur \texorpdfstring{$\rightarrow$}{vers} backend)}

\begin{table}[H]
\centering
\caption{Commandes acceptées sur \code{/ws}.}
\label{tab:commandes}
\footnotesize
\begin{tabularx}{\textwidth}{@{}l X@{}}
\toprule
\textbf{Action} & \textbf{Message} \\
\midrule
Appliquer un scénario & \code{\{"cmd":"preset","name":"controle\_auto"\}} \\
Temps de cycle d'un poste & \code{\{"cmd":"set\_time","station":2,"value":8\}} (1 à 60~s) \\
Machines en parallèle & \code{\{"cmd":"set\_servers","station":2,"value":2\}} (1 à 3) \\
Convoyeur adaptatif & \code{\{"cmd":"adaptive","value":true\}} \\
Pannes aléatoires & \code{\{"cmd":"failures","value":false\}} \\
Provoquer / réparer & \code{\{"cmd":"break","station":1\}}, \code{\{"cmd":"repair","station":1\}} \\
Vitesse, pause, remise à zéro & \code{\{"cmd":"speed","value":10\}}, \code{\{"cmd":"pause","value":true\}}, \code{\{"cmd":"reset"\}} \\
Comparer les scénarios & \code{\{"cmd":"compare","duration":3600\}} $\rightarrow$ réponse \code{compare} \\
\bottomrule
\end{tabularx}
\end{table}

\section{Messages backend \texorpdfstring{$\rightarrow$}{vers} clients (page web et scène Unity)}

\begin{itemize}
  \item \code{hello} (à la connexion)~: liste des postes et des scénarios disponibles.
  \item \code{state} (10~Hz)~: \code{config}, \code{stations} (état, file, utilisation, disponibilité, attente, \ac{MTBF}/\ac{MTTR} observés, machines), \code{products} (poste, position \code{q}/\code{p}/\code{t}, progression), \code{kpi}, \code{history} (3~dernières minutes), \code{events}, \code{exits} (produits sortis ou rejetés depuis le message précédent).
  \item \code{compare}~: \code{status} (\code{running} puis \code{done}) et résultats par scénario.
  \item \code{error}~: commande invalide.
\end{itemize}

\section{Points d'accès REST}

\begin{table}[H]
\centering
\caption{\ac{API} \ac{REST} de l'historique.}
\label{tab:rest}
\begin{tabularx}{\textwidth}{@{}l X@{}}
\toprule
\textbf{Route} & \textbf{Réponse} \\
\midrule
\code{GET /api/state} & Instantané de l'état courant (\ac{JSON}) \\
\code{GET /api/runs} & Liste des exécutions avec nombre de relevés et production \\
\code{GET /api/history?run=N} & Relevés par minute et événements de l'exécution \code{N} \\
\code{GET /api/export.csv[?run=N]} & Export \ac{CSV} (séparateur \code{;}) d'une exécution ou de tout l'historique \\
\bottomrule
\end{tabularx}
\end{table}
